\begin{theorem}[单一幂散度的渐近统计完备性]\label{thm:conclusion-binfold-single-powerdiv-asymptotic-statistical-completeness}
\leanverified{paper\_conclusion\_binfold\_single\_powerdiv\_asymptotic\_statistical\_completeness}
在 bin-fold 主线中，固定任意 \(\alpha>1\)，并记
\[
f_\alpha(t):=t^\alpha-1-\alpha(t-1).
\]
设两组稳定层输出分布 \((\mu_m)\)、\((\mu'_m)\) 分别相对于其稳定层均匀基线 \(u_m,u'_m\) 满足
\[
\lim_{m\to\infty}D_{f_\alpha}(\mu_m\Vert u_m)
=
\lim_{m\to\infty}D_{f_\alpha}(\mu'_m\Vert u'_m).
\]
则二者具有同一黄金参数 \(\varphi\)，并因而同时具有：
\[
\text{同一全 Rényi 极限族 }\{D_\beta^\infty\}_{\beta>0},
\qquad
\text{同一全 Csisz\'ar 极限锥 }\{D_f^\infty\},
\]
\[
\text{同一两点残差律 }Z^{(q)}\in\{1,\varphi^{-1}\},
\qquad
\text{同一 escort Fisher 几何},
\qquad
\text{同一涨落--响应律}.
\]
更强地，这两组实验在渐近 Le Cam 意义下属于同一统计实验。
\end{theorem}

\begin{proof}
由定理 \ref{thm:conclusion-powerdiv-scalar-recovers-escort-thermo-ladder}，任意一个固定 \(\alpha>1\) 上的极限标量
\[
D_{f_\alpha}^\infty
\]
已足以唯一恢复黄金参数 \(\varphi\)。故一旦两组分布满足同一极限幂散度，它们便具有相同的 \(\varphi\)。

随后，命题 \ref{prop:conclusion-binfold-escort-csiszar-blackwell-phi}、定理 \ref{thm:conclusion-powerdiv-scalar-recovers-escort-thermo-ladder} 与定理 \ref{thm:conclusion-binfold-escort-bayes-risk-collapse} 共同说明：全 Rényi 极限族、全 Csisz\'ar 极限锥、两点残差律、escort Fisher 几何以及涨落--响应常数，全部都是 \(\varphi\) 的显式函数。故这些极限对象在两组实验间完全一致。

最后，由定理 \ref{thm:conclusion-binfold-escort-blackwell-equivalence}，两组 bin-fold 实验族都以指数精度塌缩到同一个由 \(\varphi\) 唯一决定的 Bernoulli 两点极限实验。因而它们分别与同一极限实验的 Le Cam 距离趋于零，进而彼此的 Le Cam 距离也趋于零。故它们属于同一渐近统计实验。证毕。
\end{proof}

\begin{corollary}[bin-fold 极限实验不存在非平凡 nuisance 变形]\label{cor:conclusion-binfold-no-nontrivial-nuisance-deformation-fixed-powerdiv}
在定理 \ref{thm:conclusion-binfold-single-powerdiv-asymptotic-statistical-completeness} 的设定下，任何保持某个固定 \(\alpha>1\) 的极限幂散度
\[
D_{f_\alpha}^\infty
\]
不变的变形，都是渐近统计上平凡的；亦即，它不可能改变任一固定温度 \(q\) 上的两点极限残差、任一 \(\beta\)-Rényi 极限、任一 Csisz\'ar 极限、或 escort Fisher 曲率。
\end{corollary}

\begin{proof}
该变形既保持某个固定 \(\alpha>1\) 的极限幂散度不变，由定理 \ref{thm:conclusion-binfold-single-powerdiv-asymptotic-statistical-completeness}，便必须保持 \(\varphi\) 不变。而上述全部极限对象都已知是 \(\varphi\) 的显式函数，故它们逐一不变。于是该变形在极限实验层面不再携带独立参数。证毕。
\end{proof}
